% AUTORIDAD DE FUENTE: Exposé IX, §8.1, líneas 2109--2154;
% índices de escaneo 425--427, folios 414--416. Las capas francesas D/C
% no presentan diferencias materiales. Dos referencias internas impresas
% inequívocamente erróneas se conservan reversiblemente mediante \sisi y
% se han comunicado al custodio raíz.

\medskip
\noindent
8. \underline{La extensión de Raynaud} $A_K^{\natural o}$ \underline{en el caso no
semiestable}, \underline{y el ind-grupo} $A_K^{\flat}$

\medskip
\noindent
8.1. Sigamos suponiendo que el rasgo $S$ es completo, pero ya no
supongamos necesariamente que $A_K$ tenga reducción semiestable. Como
de costumbre, sea $K'$ una extensión de Galois finita de $K$ tal que
$A_{K'}$ tenga reducción semiestable con respecto a la normalización
$S'$ de $S$ en $K'$. Consideremos la extensión de Raynaud
$A_{K'}^{\natural}$ asociada a $A_{K'}$ (7.6.1). Es claro que el grupo
de Galois de $K'/K$ actúa por transporte de estructura sobre este
esquema en grupos; por descenso de Galois, da entonces un esquema en
grupos sobre $K$, denotado $A_{K,K'}^{\natural}$, que depende de la
elección de la extensión $K'$. Sin embargo, su componente neutra es
independiente de esta elección salvo isomorfismo canónico, por 7.6.3, y
se denotará simplemente por $A_K^{\natural o}$. El descenso demuestra
que $A_K^{\natural o}$ es una extensión

\begin{equation}\tag{8.1.1}
0 \longrightarrow T_K \longrightarrow A_K^{\natural o}
 \longrightarrow B_K \longrightarrow 0 \quad ,
\end{equation}
donde $T_K$ es un toro y $B_K$ un esquema abeliano sobre $K$. Del mismo
modo, por la definición de las partes esencialmente fija y esencialmente
tórica de $T_{\ell}(A_K)$ (4.1 y \sisi{6.5}{6.3}), los isomorfismos
(7.6.2) dan isomorfismos canónicos

\begin{equation}\tag{8.1.2}
T_{\ell}(A_K^{\natural}) \simeq T_{\ell}(A_K)^{ef} \quad ,
\qquad T_{\ell}(T_K) \simeq T_{\ell}(A_K)^{et} \quad ,
\end{equation}
y por tanto, pasando al cociente,

\begin{equation}\tag{8.1.3}
T_{\ell}(B_K) \simeq
T_{\ell}(A_K)^{ef}/T_{\ell}(A_K)^{et} \quad .
\end{equation}

%%%%%%%%%%  scan idx 426  |  folio 415  %%%%%%%%%%

Esto determina, por tanto, dos de los cocientes que aparecen en la
filtración

\begin{equation}\tag{8.1.4}
T_\ell(A_K) \supset T_\ell(A_K)^{ef}
 \supset T_\ell(A_K)^{et} \supset 0.
\end{equation}
Así, la extensión de esquemas en grupos (8.1.1) describe dos cocientes
de esta filtración. Para el tercer cociente
$T_\ell(A_K)/T_\ell(A_K)^{ef}$, $\sisi{(6.5.2)}{(6.3.2)}$ da un
isomorfismo canónico

\begin{equation}\tag{8.1.5}
T_\ell(A_K)/T_\ell(A_K)^{ef}
 \simeq D(T_\ell(A'_K)^{et}) \quad ,
\end{equation}
donde $D=\underline{\operatorname{Hom}}(-,T_\ell(\mathbb{G}_m))$
denota la dualidad de Cartier (en el sentido de los progrupos de
Barsotti--Tate), y $A'_K$ denota el esquema abeliano dual de $A_K$.
Consideremos entonces la sucesión exacta de tipo (8.1.1)

\begin{equation}\tag{8.1.6}
0 \longrightarrow T'_K \longrightarrow {A'_K}^{\natural o}
 \longrightarrow B'_K \longrightarrow 0
\end{equation}
asociada al esquema abeliano $A'_K$ sobre $K$, junto con los
isomorfismos correspondientes de tipo (8.1.2). La fórmula (8.1.5) da
entonces

\begin{equation}\tag{8.1.7}
T_\ell(A_K)/T_\ell(A_K)^{ef}
 \simeq D(T_\ell(T'_K))
 \simeq \underline{M}_K\otimes_{\mathbb{Z}}\mathbb{Z}_\ell \quad ,
\end{equation}
donde, generalizando la notación introducida en (5.5.2), hemos puesto

\begin{equation}\tag{8.1.8}
\underline{M}_K=D(T'_K) \quad , \qquad
\underline{M}'_K=D(T_K) \quad ,
\end{equation}
y donde $D$ denota ahora la dualidad de Cartier
$\underline{\operatorname{Hom}}(-,\mathbb{G}_{mK})$ en el sentido de
los esquemas en grupos ordinarios. Así, $\underline{M}_K$ y
$\underline{M}'_K$ son grupos constantes torcidos sobre $K$ que se
vuelven constantes sobre $K'$ e isomorfos a grupos de la forma
$\mathbb{Z}^{r}$, donde $r=\dim T=\dim T'$.
